\documentclass[11pt]{article}
\usepackage[a4paper,margin=27mm]{geometry}
\usepackage{amsmath,amssymb,amsthm}
\usepackage{microtype}
\usepackage[hidelinks]{hyperref}
\newtheorem{theorem}{Theorem}
\newtheorem{lemma}[theorem]{Lemma}
\newcommand{\PP}{\mathbf P}
\newcommand{\OO}{\mathcal O}
\newcommand{\mult}{\operatorname{mult}}
\newcommand{\Hom}{\operatorname{Hom}}
\newcommand{\Aut}{\operatorname{Aut}}
\newcommand{\Spec}{\operatorname{Spec}}
\title{Positive tangent-bundle Seshadri constants\protect\\characterize projective space}
\author{Alec Kriebel\\[2pt]\small Independent researcher\\
\small\href{https://orcid.org/0009-0001-9320-500X}{ORCID: 0009-0001-9320-500X}}
\date{12 October 2026\quad (version 1.0)}
\begin{document}
\maketitle
\begin{abstract}
Let $X$ be a smooth integral projective variety over an algebraically closed
field of arbitrary characteristic. We prove that positivity of the relative
Seshadri constant of $T_X$ at one closed point implies $X\simeq\PP^{\dim X}$,
resolving the conjecture of Fulger and Murayama. An absolute minimum of ample
degree among rational curves through the specified point gives a proper
pointed family without boundary. Its very-free members yield surjective
evaluation. A contracted marking section and line-bundle descent force
Picard number one and ampleness of $-K_X$, reducing the problem to the
established Fano case. No general-point or characteristic-zero assumption is
used. This is an unrefereed preprint developed and checked extensively with
AI tools.
\end{abstract}

\section{The statement and its context}

For a vector bundle $E$ on a projective variety $X$, let
$\rho:\PP(E)=\operatorname{Proj}_X\operatorname{Sym} E\to X$ use the
quotient convention, and write $\xi=c_1(\OO_{\PP(E)}(1))$. The relative
Seshadri constant of Fulger--Murayama is
\begin{equation}\label{eq:definition}
 \epsilon(E;x)=\inf_{\Gamma}\frac{\xi\cdot\Gamma}
 {\mult_x(\rho_*\Gamma)},
\end{equation}
where $\Gamma$ ranges over integral curves in $\PP(E)$ whose images contain
$x$ and which are not contracted by $\rho$. If $\Gamma\to C$ has total
generic degree $q$, then $\rho_*\Gamma=q[C]$; the degree includes its
inseparable factor. Global nefness of $E$ is not required.

\begin{theorem}\label{thm:main}
Let $X$ be a smooth integral projective variety of dimension $n\geq1$ over
an algebraically closed field $k$ of arbitrary characteristic. If
$\epsilon(T_X;x)>0$ at a closed point $x\in X$, then $X\simeq_k\PP^n$.
\end{theorem}

This is \cite[Conjecture 4.9]{FM}, also stated in Fulger's contribution to
the 2019 Oberwolfach workshop \cite[p.~3289]{OWR}. Fulger--Murayama prove
the Fano case in every characteristic, the surface case, and the
general-point case in characteristic zero. Chang proves the smooth
projective toric case \cite[Theorem 0.2]{Chang}. The conclusion here retains
the original arbitrary point and removes the additional Fano hypothesis.

Proper pointed rational-curve families and their numerical-generation
properties are standard tools; see, for example,
\cite[Proposition 2.20(b)]{AO}. We give the numerical argument explicitly
using a genuine contracted section. The proof combines these tools with a
minimum taken over \emph{all} rational images through the given point and
a characteristic-free construction of the actual marked family.
The final characterization is imported from the established Fano clause
\cite[Proposition 4.8(1)]{FM}, which uses Mori's theorem as presented in
\cite[Theorem V.3.2]{Kollar}. The conjecture itself is never an input.

\section{A complete family through the specified point}

\begin{lemma}\label{lem:free}
Every normalization $f:\PP^1\to C\hookrightarrow X$ of an integral
rational curve $C$ through $x$ has $f^*T_X$ ample.
\end{lemma}
\begin{proof}
Write $f^*T_X=\bigoplus_{i=1}^n\OO_{\PP^1}(a_i)$, and put
$m=\mult_x C$. A quotient onto $\OO(a_i)$ defines a lift
$s_i:\PP^1\to\PP(T_X)$ of $f$. Its image $\Gamma_i$ is birationally
parametrized by $s_i$: its function field lies between the equal fields
$k(C)$ and $k(\PP^1)$. Hence
$\xi\cdot\Gamma_i=a_i$ and $\rho_*\Gamma_i=[C]$.
Equation~\eqref{eq:definition} gives
$a_i/m\geq\epsilon(T_X;x)>0$. Thus each integer $a_i\geq1$.
This argument works in every characteristic, including singular rational
images, without replacing Frobenius-normalized slopes on general curves
by ordinary slopes.
\end{proof}

The proof of \cite[Corollary 4.6]{FM} produces an integral rational curve
through the \emph{same} point $x$ from the hypothesis of
Theorem~\ref{thm:main}. Fix a very ample line bundle $A$ on $X$, and set
\begin{equation}\label{eq:minimum}
d=\min\{A\cdot C:C\subset X\text{ is integral rational and }x\in C\}.
\end{equation}
This is an attained positive integer.

\begin{lemma}\label{lem:boundary}
Every geometric genus-zero stable map of $A$-degree $d$ with one marking
evaluating to $x$ has smooth domain $\PP^1$ and is birational onto its image.
\end{lemma}
\begin{proof}
First the minimum in \eqref{eq:minimum} survives any algebraically closed
extension $K/k$. A rational image through $x_K$ of degree $c<d$ would give
a degree-$c$ normalization map $\PP^1_K\to X_K$, reparametrized so that
$0\mapsto x_K$. The degree-$c$ pointed Hom scheme is of finite type over
$k$. Its nonemptiness would give a $k$-point, since $k$ is algebraically
closed. The image of that map is a rational curve through $x$ of degree
at most $c$, contradicting \eqref{eq:minimum}. Rationality of the image
follows from L\"uroth's theorem, also in positive characteristic.

Consider a geometric map $(g:D\to X,p)$ as in the statement. Its dual
graph is a tree of projective lines. Starting at the marked component,
follow components contracted to $x$ until reaching a nonconstant component
$D_v$. Its image $C_v$ is a rational curve through $x$. If $e_w$ is the
full generic degree of each nonconstant component onto its image, then
\[
d=\sum_{w\text{ nonconstant}} e_w(A\cdot C_w)
 \geq e_v(A\cdot C_v)\geq A\cdot C_v\geq d.
\]
All summands are positive integers. There is therefore just one
nonconstant component, and it maps birationally to a degree-$d$ image.
Any connected contracted subtree with $t>0$ vertices attaches to it along
one edge and has at most one marking. Its total number of special-point
incidences is at most $2(t-1)+1+1=2t$, whereas stability requires at least
$3t$. Such a subtree cannot exist. This includes a marking initially on
a contracted component and excludes inseparable multiple covers by their
full degree.
\end{proof}

\begin{lemma}\label{lem:family}
There are an integral projective scheme $H$, a smooth projective-line
bundle $\pi:U\to H$, a section $\sigma:H\to U$, and a surjective
morphism $e:U\to X$ such that $e\sigma=x$. Every fiber map is birational
onto a degree-$d$ rational curve through $x$.
\end{lemma}
\begin{proof}
Abramovich--Oort \cite[Theorem 2.8 and \S2.5]{AbramovichOort} give a proper
finite-presentation Artin stack of one-marked genus-zero maps to $X$ of
$A$-degree $d$, with a projective coarse scheme $Q$. Let $\mathcal M_x$
be the actual evaluation fiber over $x$. It is proper and nonempty.
Lemma~\ref{lem:boundary} describes all its geometric objects.

Their stabilizer group schemes are trivial. A geometric automorphism
preserving a birational map is the identity on the open set where the map
is an isomorphism onto its image. Infinitesimal automorphisms are vector
fields on $\PP^1$ vanishing at the mark and killed by $dg$. The
normalization map has generically injective differential, so such a
vector field is zero. A finite group scheme with one geometric point and
zero tangent space is the identity: its local maximal ideal satisfies
$\mathfrak m/\mathfrak m^2=0$, and Nakayama applies.
This also makes inertia trivial over nonreduced bases. On an atlas its
finite algebra has an augmentation split by the identity section as
$\OO\oplus J$. The finite module $J$ has zero geometric fibers and
therefore vanishes by Nakayama.

Trivial inertia makes $\mathcal M_x$ an algebraic space
\cite[Tag 04SZ]{Stacks}. Its natural morphism to $Q$ is proper, because
$\mathcal M_x$ is proper over $k$ and $Q$ is separated. It is quasi-finite:
coarse geometric points identify isomorphism classes, so each geometric
fiber contains at most one point. Tags 03XX and 02LS of \cite{Stacks}
make $\mathcal M_x$ a scheme finite over $Q$, hence projective.
No commutation of coarse moduli with a closed evaluation fiber in wild
characteristic is assumed.

The stack's universal curve restricts to an actual smooth proper
$\PP^1$-family on $\mathcal M_x$, with marking section. The marking is
a relative effective Cartier divisor. Its line bundle has fiber degree
one; cohomology and base change give a rank-two pushforward and a
fiberwise surjective evaluation map. This identifies the family with
the quotient projective bundle of that pushforward, also over nonreduced
bases. One can check the identification by noting that the induced
morphism is proper and fiberwise an isomorphism, hence finite; Nakayama
makes its algebra map surjective, and flatness over the base makes its
defining ideal have zero fibers, so that ideal vanishes.

Choose a degree-$d$ normalization $f:\PP^1\to X$ with $f(0)=x$.
Let $R_d$ be the degree-$d$ Hom scheme with $0\mapsto x$.
Its morphism to $\mathcal M_x$ is the torsor of pointed parametrizations
under the smooth group $\Aut(\PP^1,0)$. Indeed the projective-line bundle
with section just constructed is locally a pointed trivial bundle.
By Lemma~\ref{lem:free},
\[
H^1(\PP^1,f^*T_X(-0))=0,\qquad
H^1(\PP^1,f^*T_X(-0-\infty))=0.
\]
Fixed-source deformation theory makes $R_d$ smooth at $f$ and evaluation
at $\infty$ smooth there. For the latter assertion the obstruction is
the second displayed group: local lifts into smooth $X$ respecting the
two prescribed values differ by a \v Cech cocycle in
$f^*T_X(-0-\infty)$, whose vanishing obstruction allows the lifts to
glue. This is infinitesimal smoothness, not generic smoothness in
characteristic zero; see also \cite[Proposition 4.7]{Gounelas}.

Smooth descent along the parametrization torsor makes $\mathcal M_x$
smooth at $[f,0]$. Let $H$ be its unique reduced irreducible component
there, and restrict the universal curve to $U\to H$. A neighborhood
of $f$ in $R_d$ maps into $H$. Its smooth evaluation at $\infty$ factors
through $e:U\to X$, so $e(U)$ contains a nonempty open subset of $X$.
It is closed because $U$ is proper; irreducibility of $X$ gives $e(U)=X$.
The total space $U$ is integral: smoothness over reduced $H$ gives
reducedness, and flatness with geometrically integral fibers over integral
$H$ gives irreducibility. It is projective because it is a projective-line
bundle over projective $H$. All asserted fiber degrees and birationality
follow from Lemma~\ref{lem:boundary}.
\end{proof}

\section{The contracted section and the Fano reduction}

We spell out the standard numerical-generation argument for the family
above. In particular it tests curves avoiding $x$, as well as curves of
positive genus.

\begin{lemma}\label{lem:numerical}
For every Cartier divisor $D$ on $X$, put $b=D\cdot C$, where $C$ is the
image of a fiber of $U\to H$. Then $b$ is independent of the fiber and
\[
 d[D]=b[A]\quad\text{in }N^1(X)_{\mathbf R}.
\]
Consequently $\rho(X)=1$.
\end{lemma}
\begin{proof}
Degrees of a line bundle on fibers of a proper flat curve family are
locally constant. They are constant here since $H$ is integral. The
line bundle $L=\OO_X(dD)\otimes A^{-b}$ thus pulls back to a bundle
$M=e^*L$ of degree zero on every geometric $\PP^1$-fiber. Each such
restriction is trivial and has $h^0=1$, $h^1=0$. Proper flat cohomology
and base change give a line bundle $B=\pi_*M$ with
$\pi^*B\simeq M$; the adjunction map is an isomorphism on every
geometric fiber and hence everywhere. The applicable perfect-base-change
result is \cite[Tag 0B91]{Stacks}; direct fiberwise descent is also
given by Tags 0EX7 and 0GKA. Pulling back by $\sigma$ gives
\[
 B\simeq\sigma^*e^*L\simeq L|_x\otimes_k\OO_H\simeq\OO_H.
\]
Thus $e^*L$ is trivial.

Let $\Gamma\subset X$ be any integral curve and let $\eta$ be its
generic point. Surjectivity of $e$ makes $U_\eta$ nonempty. A closed
point of this finite-type $k(\Gamma)$-scheme has residue field $K$ finite
over $k(\Gamma)$, of total degree $q>0$. Its reduced closure in $U$ is
a proper integral curve $\Gamma'$ dominating $\Gamma$ with generic
degree $q$. The degree projection formula gives
\[
 0=\deg(e^*L|_{\Gamma'})=q\deg(L|_\Gamma).
\]
These are integer degrees, so the conclusion remains valid even if the
characteristic divides $q$. Hence $L$ is numerically trivial on $X$,
proving the asserted relation. The ample class $[A]$ is nonzero since
$A\cdot C=d>0$, so the Picard number is exactly one.
\end{proof}

\begin{proof}[Proof of Theorem~\ref{thm:main}]
Use Lemma~\ref{lem:numerical} with $D=-K_X$. By
Lemma~\ref{lem:free}, its degree on the original rational curve is
$b=\deg f^*T_X=\sum_i a_i>0$. Thus
$-K_X\equiv cA$ with $c=b/d>0$.
For every integral subvariety $Z\subset X$ of dimension $r\geq1$,
numerical equivalence gives
\[
 (-K_X)^r\cdot Z=c^r A^r\cdot Z>0.
\]
For completeness, subtracting the two powers expresses their difference
as $(-K_X-cA)$ intersected with
$\sum_{j=0}^{r-1}(-K_X)^{r-1-j}(cA)^j\cap[Z]$, a real one-cycle class;
the difference therefore has degree zero. Since $X$ is smooth,
$-K_X$ is Cartier. The Nakai--Moishezon criterion, valid in every
characteristic (see \cite[Theorem 1.3]{FujinoMiyamoto}), makes $-K_X$
ample. The established Fano clause \cite[Proposition 4.8(1)]{FM} now
gives $X\simeq_k\PP^n$ at the original point and over the original field.
\end{proof}

When $n=1$, one can instead use
$\epsilon(T_X;x)=\deg T_X=2-2g$, so positivity forces $g=0$.
Conversely projective space has positive tangent Seshadri constant:
it is $2$ for $\PP^1$ and $1$ for $\PP^n$, $n\geq2$
\cite[Example 4.2]{FM}. Smoothness is essential; the singular toric
examples in \cite[Example 0.4]{Chang} do not satisfy this theorem's
hypotheses.

\section*{Verification, priority, and AI-use disclosure}

The supporting package contains the source, the original frozen candidate,
reproducible integrity and finite controls, independent mathematical audit
reports, and a dated priority-search ledger. The finite controls supplement
the written proof; they do not formally verify the geometric theorem.
A bounded primary-source priority audit on 12 October 2026 identified no
earlier full resolution or explicitly matching immediate consequence. It
does not certify first historical priority. The standard numerical-generation
principle and the previously established special cases are credited above.
The full Koll\'ar monograph and full CMSB text were not independently obtained
in this audit; the exact published Fulger--Murayama propositions used here
were directly inspected. An in-preparation reference on moving Seshadri
constants had no inspected public contents.

AI tools were used extensively in developing the proof, checking references
and calculations, performing internal adversarial reviews, and preparing the
manuscript and verification package. Responsibility for the claims and
citations rests with the author. This preprint is unrefereed; internal
AI-assisted review is not external human peer review, and no formal
machine-checked proof is asserted.

\begin{thebibliography}{9}
\bibitem{FM} M. Fulger and T. Murayama,
\emph{Seshadri constants for vector bundles}, J. Pure Appl. Algebra
\textbf{225} (2021), no.~4, 106559.
\href{https://doi.org/10.1016/j.jpaa.2020.106559}{doi:10.1016/j.jpaa.2020.106559}.
Final-numbering public manuscript:
\url{https://par.nsf.gov/servlets/purl/10198624}.
\bibitem{OWR} M. Fulger (joint with T. Murayama),
\emph{Seshadri constants for vector bundles}, in
\emph{Mini-Workshop: Seshadri Constants}, Oberwolfach Reports
\textbf{16} (2019), Report 53, 3287--3289.
\href{https://doi.org/10.4171/OWR/2019/53}{doi:10.4171/OWR/2019/53}.
\bibitem{Chang} C.-W. Chang,
\emph{The Seshadri constants of tangent sheaves on toric varieties},
J. Pure Appl. Algebra \textbf{229} (2025), no.~9, 108046.
\href{https://doi.org/10.1016/j.jpaa.2025.108046}{doi:10.1016/j.jpaa.2025.108046};
\href{https://arxiv.org/abs/2211.17172}{arXiv:2211.17172v2}.
\bibitem{AO} M. Andreatta and G. Occhetta,
\emph{Fano manifolds with long extremal rays}, Asian J. Math.
\textbf{9} (2005), no.~4, 523--543.
\href{https://doi.org/10.4310/AJM.2005.v9.n4.a5}{doi:10.4310/AJM.2005.v9.n4.a5}.
\bibitem{AbramovichOort} D. Abramovich and F. Oort,
\emph{Stable maps and Hurwitz schemes in mixed characteristic},
\href{https://arxiv.org/abs/math/9808074}{arXiv:math/9808074} (1998),
Theorem 2.8 and \S2.5.
\bibitem{Kollar} J. Koll\'ar,
\emph{Rational Curves on Algebraic Varieties}, Ergebnisse der Mathematik
und ihrer Grenzgebiete, 3. Folge, vol.~32, Springer, 1996.
\href{https://doi.org/10.1007/978-3-662-03276-3}{doi:10.1007/978-3-662-03276-3}.
\bibitem{Gounelas} F. Gounelas,
\emph{Free curves on varieties}, Documenta Math. \textbf{21} (2016),
287--308, Proposition 4.7.
\url{https://ems.press/content/serial-article-files/26314?nt=1}.
\bibitem{FujinoMiyamoto} O. Fujino and A. Miyamoto,
\emph{Nakai--Moishezon ampleness criterion for real line bundles},
Theorem 1.3, author manuscript.
\url{https://www.math.kyoto-u.ac.jp/~fujino/Nakai-Moishezon3.pdf}.
\bibitem{Stacks} The Stacks Project Authors, \emph{The Stacks Project},
Tags 04SZ, 03XX, 02LS, 0B91, 0EX7, and 0GKA.
\url{https://stacks.math.columbia.edu/}.
\end{thebibliography}
\end{document}
